% ============================================================
% SECCIÓN 4
% ============================================================
\section{Espectro de Diseño según NCh433}
\label{sec:espectro}

El espectro de diseño se construye conforme al artículo 6.3 de la norma 
NCh433~\cite{NCh433:1996}, utilizando los parámetros de la 
Tabla~\ref{tab:parametros_sismicos}.

\subsection{Espectro Elástico}

La aceleración espectral elástica $S_a(T)$ se define como:

\begin{equation}
    S_a(T) = A_0 \cdot I \cdot \alpha(T)
\end{equation}

donde $\alpha(T)$ es el coeficiente de amplificación espectral dado por:

\begin{equation}
    \alpha(T) = \frac{1+4.5 \left(\frac{T_n}{T_0} \right)^p}{1 +\left( \frac{T_n}{T_0}\right)^3}
\end{equation}


\subsection{Espectro Reducido}

El espectro reducido se obtiene dividiendo el espectro elástico por el factor 
de corrección $R^*$ calculado en la Sección~\ref{sec:factores R}, siendo distinto 
para cada dirección de análisis:

\begin{equation}
    S_{a,\text{red}}(T) = \frac{S_a(T)}{R^*}
\end{equation}

Los factores de reducción definitivos, ajustados para cumplir con el corte basal 
mínimo exigido por la norma, son:

\begin{equation}
    R^*_X = 2{,}22 \qquad R^*_Y = 2{,}69
\end{equation}

En la Figura~\ref{fig:espectros} se presentan el espectro elástico y los espectros 
reducidos para ambas direcciones. Se indican además los períodos fundamentales de 
traslación ($T_X = 1{,}14\ \text{s}$, $T_Y = 0{,}853\ \text{s}$), 
utilizados en el cálculo de $R^*$.

\begin{figure}[H]
    \centering
    \includegraphics[width=0.85\textwidth]{IMG/Espectro_NCh433_Grupo1.png}
    \caption{Espectro elástico y espectros reducidos para las direcciones X e Y, 
             según NCh433~\cite{NCh433:1996}.}
    \label{fig:espectros}
            \vspace{0.001cm}
   
    \text{Fuente: Elaboración propia}
\end{figure}